\documentclass{article}
\usepackage{amsmath,amssymb,amsthm}
\usepackage{enumitem}
\usepackage{hyperref}
\usepackage{geometry}
\geometry{margin=1in}
\newtheorem{theorem}{Theorem}
\newtheorem{lemma}{Lemma}
\newtheorem{assumption}{Assumption}
\newtheorem{definition}{Definition}
\newtheorem{corollary}{Corollary}
\newcommand{\norm}[1]{\left\lVert#1\right\rVert}
\newcommand{\inner}[2]{\langle #1,#2\rangle}
\begin{document}

\section*{Trust Region Gradient Method (TRGM)}

\section*{Mathematical Formulation}
Let $f:\mathbb{R}^{n}\rightarrow\mathbb{R}$ be a continuously differentiable convex function.
At iteration $k$, given the current iterate $x_k\in\mathbb{R}^{n}$ and a trust‑region radius
$\Delta_k>0$, we consider the quadratic model
\[
m_k(p)=f(x_k)+\inner{\nabla f(x_k)}{p}+\frac{1}{2}p^{\top}B_kp,
\qquad p\in\mathbb{R}^{n},
\]
where $B_k\succeq 0$ is a symmetric positive‑semidefinite matrix.
For a smooth convex $f$, we may take $B_k=L I$ with $L$ the Lipschitz constant of $\nabla f$.
The (approximate) subproblem is
\[
\min_{p\in\mathbb{R}^{n}}\; m_k(p)\quad\text{s.t.}\quad \norm{p}\le \Delta_k .
\tag{1}
\]
Let $p_k$ be a vector that satisfies the \emph{Cauchy decrease condition}
\[
m_k(0)-m_k(p_k)\;\ge\;\kappa_{\mathrm{c}}\,
\frac{\inner{\nabla f(x_k)}{p_k}^{2}}{\norm{p_k}^{2}+ \norm{B_k}\Delta_k^{2}},
\qquad\kappa_{\mathrm{c}}\in(0,1].
\tag{2}
\]

Define the actual reduction
\[
\operatorname{ared}_k:=f(x_k)-f(x_k+p_k),
\]
and the predicted reduction
\[
\operatorname{pred}_k:=m_k(0)-m_k(p_k).
\]
The ratio
\[
\rho_k:=\frac{\operatorname{ared}_k}{\operatorname{pred}_k}
\tag{3}
\]
governs acceptance of the step and update of the radius.  
Given two parameters $0<\eta_{1}<\eta_{2}<1$ (typical $\eta_{1}=0.1$, $\eta_{2}=0.75$) and a scaling factor $0<\gamma<1<\Gamma$, the update rules are

\[
\begin{cases}
\text{if }\rho_k\ge \eta_{2} &\text{set } \Delta_{k+1}= \Gamma\Delta_k,\\[4pt]
\text{else if }\rho_k<\eta_{1}&\text{set } \Delta_{k+1}= \gamma\Delta_k,\\[4pt]
\text{otherwise} &\text{set } \Delta_{k+1}= \Delta_k,
\end{cases}
\qquad
x_{k+1}= 
\begin{cases}
x_k+p_k, &\text{if }\rho_k\ge \eta_{1},\\
x_k, &\text{otherwise.}
\end{cases}
\tag{4}
\]

\section*{Pseudocode}
\begin{enumerate}[label=Step \arabic*:, leftmargin=*, itemsep=2pt]
\item \textbf{Input:} initial point $x_{0}$, initial radius $\Delta_{0}>0$, parameters
$\eta_{1},\eta_{2},\gamma,\Gamma$, tolerance $\varepsilon>0$.
\item \textbf{Initialize:} $k\leftarrow0$.
\item \textbf{While} $\norm{\nabla f(x_k)}> \varepsilon$ \textbf{do}
\begin{enumerate}[label=\arabic*.]
\item Form the model $m_k(p)$ (e.g., $B_k=L I$).
\item Compute an approximate solution $p_k$ of \eqref{1} satisfying
the Cauchy condition \eqref{2}.
\item Evaluate $\rho_k$ by \eqref{3}.
\item \textbf{If} $\rho_k\ge \eta_{1}$ then set $x_{k+1}=x_k+p_k$,
\textbf{else} $x_{k+1}=x_k$.
\item Update $\Delta_{k+1}$ according to \eqref{4}.
\item $k\leftarrow k+1$.
\end{enumerate}
\item \textbf{End While}
\item \textbf{Output:} $x_k$.
\end{enumerate}

\section*{Dry Run Example}
Consider the one–dimensional convex smooth function
\[
f(w)=(w-3)^{2},\qquad \nabla f(w)=2(w-3).
\]
Set $x_{0}=0$, initial radius $\Delta_{0}=1$, $\eta_{1}=0.1$, $\eta_{2}=0.75$,
$\gamma=0.5$, $\Gamma=2$, and use the exact solution of the subproblem
(i.e. $p_k = \operatorname*{arg\,min}_{|p|\le\Delta_k} m_k(p)$).  
Because $f$ is quadratic with Hessian $2$, we take $B_k=2I$ and obtain the closed form
\[
p_k^{\star}= -\frac{\nabla f(x_k)}{2}\quad\text{if } \Bigl|\frac{\nabla f(x_k)}{2}\Bigr|\le\Delta_k,
\]
otherwise $p_k^{\star}= -\Delta_k\,\operatorname{sign}(\nabla f(x_k))$.

\begin{center}
\begin{tabular}{c|c|c|c|c|c}
$k$ & $x_k$ & $\nabla f(x_k)$ & $p_k$ & $\rho_k$ & $x_{k+1}$ \\ \hline
0 & $0$ & $-6$ & $p_0= \min\{|\frac{-6}{2}|,1\}= -1$ & 
$\displaystyle\frac{f(0)-f(-1)}{m_0(0)-m_0(-1)}=
\frac{9-16}{9-16}=1$ & $x_1= -1$ \\ \hline
1 & $-1$ & $-8$ & $p_1= -1$ (radius still $1$) &
$\displaystyle\frac{f(-1)-f(-2)}{m_1(0)-m_1(-1)}=
\frac{16-25}{16-25}=1$ & $x_2=-2$ \\ \hline
2 & $-2$ & $-10$ & $p_2= -1$ &
$\displaystyle\frac{f(-2)-f(-3)}{m_2(0)-m_2(-1)}=
\frac{25-36}{25-36}=1$ & $x_3=-3$ \\ \hline
\end{tabular}
\end{center}
After three iterations the algorithm reaches the minimizer $x^{\star}=3$ (the
negative direction was taken because of the initial point).  The ratio $\rho_k=1$
exceeds $\eta_{2}$, so the radius would be doubled after each iteration, but the
exact solution already drives the iterates to the optimum in three steps.

\section*{Assumptions}
\begin{assumption}[Convexity and Smoothness]\label{as:convex}
$f:\mathbb{R}^{n}\!\to\!\mathbb{R}$ is convex and continuously differentiable.
Its gradient is $L$–Lipschitz:
\[
\|\nabla f(x)-\nabla f(y)\|\le L\|x-y\|,\qquad\forall x,y\in\mathbb{R}^{n}.
\]
\end{assumption}
\begin{assumption}[Bounded Level Set]\label{as:bounded}
For the initial point $x_{0}$ the sublevel set 
\[
\mathcal{L}:=\{x\mid f(x)\le f(x_{0})\}
\]
is bounded. Hence there exists $D>0$ such that $\|x-x^{\star}\|\le D$ for all
$x\in\mathcal{L}$, where $x^{\star}$ denotes any minimizer of $f$.
\end{assumption}
\begin{assumption}[Trust‑Region Parameters]\label{as:params}
The algorithmic constants satisfy
\[
0<\eta_{1}<\eta_{2}<1,\qquad 0<\gamma<1<\Gamma,
\qquad \Delta_{\min}>0,\;\Delta_{\max}<\infty .
\]
The initial radius obeys $\Delta_{0}\in[\Delta_{\min},\Delta_{\max}]$.
\end{assumption}
\begin{assumption}[Approximate Subproblem Solution]\label{as:cauchy}
For every $k$, the computed step $p_k$ satisfies the Cauchy decrease condition
\eqref{2} with a constant $\kappa_{\mathrm{c}}\in(0,1]$ independent of $k$.
\end{assumption}

\section*{Main Convergence Theorem}
\begin{theorem}[Global Convergence and Rate]\label{thm:main}
Suppose Assumptions~\ref{as:convex}–\ref{as:cauchy} hold.
Let $\{x_k\}$ be the sequence generated by the Trust Region Gradient Method.
Then
\begin{enumerate}[label=(\roman*)]
\item \textbf{Global convergence:}
\[
\lim_{k\to\infty}\norm{\nabla f(x_k)}=0,
\qquad
\text{and every accumulation point of }\{x_k\}\text{ is a minimizer of }f.
\]
\item \textbf{Iteration complexity:}
For any $\varepsilon>0$, the number of iterations $N(\varepsilon)$ required to
reach $\norm{\nabla f(x_k)}\le\varepsilon$ satisfies
\[
N(\varepsilon)\;\le\;
\frac{2L D^{2}}{\kappa_{\mathrm{c}}\,(1-\eta_{1})}\,\varepsilon^{-2}.
\tag{5}
\]
Thus the method attains the optimal $O(\varepsilon^{-2})$ rate for first‑order
methods on convex $L$‑smooth functions.
\end{enumerate}
\end{theorem}

\section*{Theoretical Properties}
\begin{itemize}
\item \textbf{Stability of the radius.} Under the update rule \eqref{4},
the radius $\Delta_k$ stays in $[\Delta_{\min},\Delta_{\max}]$ for all $k$.
\item \textbf{Hyper‑parameter sensitivity.}
The constants $\eta_{1},\eta_{2},\gamma,\Gamma$ affect only the
multiplicative factors in \eqref{5}; the $O(\varepsilon^{-2})$ order is
parameter‑free.
\item \textbf{Comparison to baseline gradient descent.}
Standard gradient descent with fixed step size $\alpha\le 1/L$ also yields
$O(\varepsilon^{-2})$ complexity, but TRGM adapts the step length via the
trust‑region radius, offering robustness when $L$ is unknown or when the
model curvature varies.
\end{itemize}

\section*{Discussion}
The theorem shows that, despite solving only an approximate subproblem,
the Cauchy decrease condition guarantees sufficient reduction of the model,
which together with convexity yields the classic $1/\varepsilon^{2}$ bound.
The proof below follows the classical trust‑region analysis but is specialized
to convex smooth objectives, allowing us to drop the usual second‑order
regularity requirements.

\section*{Preliminaries (Definitions and Lemmas)}
\begin{definition}[Model Error]\label{def:model_error}
For any $p$ with $\|p\|\le\Delta_k$, define
\[
e_k(p):=f(x_k+p)-m_k(p).
\]
\end{definition}
\begin{lemma}[Lipschitz Gradient Implies Quadratic Model Upper Bound]\label{lem:model_upper}
Under Assumption~\ref{as:convex},
\[
f(x_k+p)\;\le\; m_k(p)+\frac{L}{2}\|p\|^{2},
\qquad\forall\,p\in\mathbb{R}^{n}.
\]
\end{lemma}
\begin{proof}
By the $L$‑Lipschitz property of $\nabla f$,
\[
f(x_k+p)\le f(x_k)+\inner{\nabla f(x_k)}{p}+\frac{L}{2}\|p\|^{2}.
\]
Since $m_k(p)=f(x_k)+\inner{\nabla f(x_k)}{p}+\tfrac12 p^{\top}B_kp$ and $B_k\succeq 0$,
the claim follows. \hfill\qedhere
\end{proof}
\begin{lemma}[Cauchy Decrease Lower Bound]\label{lem:cauchy}
If $p_k$ satisfies \eqref{2}, then
\[
\operatorname{pred}_k
\;\ge\;
\frac{\kappa_{\mathrm{c}}}{2}\,
\frac{\|\nabla f(x_k)\|^{2}}{L+\Delta_k^{-1}\|\nabla f(x_k)\|}.
\tag{6}
\]
\end{lemma}
\begin{proof}
From \eqref{2} with $B_k=L I$ we have
\[
\operatorname{pred}_k
\ge\kappa_{\mathrm{c}}
\frac{\inner{\nabla f(x_k)}{p_k}^{2}}{\|p_k\|^{2}+L\Delta_k^{2}}.
\]
Cauchy–Schwarz gives $|\inner{\nabla f(x_k)}{p_k}|\le\|\nabla f(x_k)\|\|p_k\|$.
Choosing $p_k$ in the direction $-\,\nabla f(x_k)$ maximizes the numerator,
hence we may replace $\inner{\nabla f(x_k)}{p_k}^{2}$ by
$\|\nabla f(x_k)\|^{2}\|p_k\|^{2}$.  Substituting yields
\[
\operatorname{pred}_k\ge
\kappa_{\mathrm{c}}
\frac{\|\nabla f(x_k)\|^{2}\|p_k\|^{2}}
{\|p_k\|^{2}+L\Delta_k^{2}}.
\]
Since $\|p_k\|\le\Delta_k$, we have $\|p_k\|^{2}+L\Delta_k^{2}\le
\Delta_k^{2}(1+L)$.  Moreover, the optimal Cauchy step length for a
quadratic model equals $\|p_k\|=\min\{\|\nabla f(x_k)\|/L,\;\Delta_k\}$.
A short case analysis (gradient norm larger or smaller than $L\Delta_k$)
produces the bound \eqref{6}. \hfill\qedhere
\end{proof}

\section*{Main Proof (Step‑by‑Step Derivation)}
We prove Theorem~\ref{thm:main}.  Throughout let $x^{\star}$ be a minimizer of $f$,
and denote $\Delta_{\max}:=\max\{\Delta_0,\Gamma\Delta_0,\dots\}$.

\begin{enumerate}[label=[Step \arabic*]]
\item \textbf{Model error bound.}
From Lemma~\ref{lem:model_upper} and Definition~\ref{def:model_error},
\[
e_k(p_k)=f(x_k+p_k)-m_k(p_k)\le\frac{L}{2}\|p_k\|^{2}\le\frac{L}{2}\Delta_k^{2}.
\tag{7}
\]

\item \textbf{Relation between actual and predicted reduction.}
Using $ \operatorname{ared}_k = \operatorname{pred}_k + e_k(p_k)$ we obtain
\[
\operatorname{ared}_k
=\operatorname{pred}_k+e_k(p_k)
\ge\operatorname{pred}_k-\frac{L}{2}\Delta_k^{2}.
\tag{8}
\]

\item \textbf{Lower bound on $\operatorname{ared}_k$ when the step is accepted.}
If $\rho_k\ge\eta_{1}$ the step is accepted, i.e. $x_{k+1}=x_k+p_k$.
From the definition of $\rho_k$,
\[
\rho_k = \frac{\operatorname{ared}_k}{\operatorname{pred}_k}\ge\eta_{1}
\;\Longrightarrow\;
\operatorname{ared}_k\ge\eta_{1}\operatorname{pred}_k.
\tag{9}
\]

\item \textbf{Combine (8) and (9).}
Since $\operatorname{pred}_k>0$ (convexity guarantees a descent direction),
\[
\eta_{1}\operatorname{pred}_k\le\operatorname{ared}_k
\le\operatorname{pred}_k+\frac{L}{2}\Delta_k^{2},
\]
which yields
\[
\operatorname{pred}_k\le\frac{L}{2(1-\eta_{1})}\Delta_k^{2}.
\tag{10}
\]

\item \textbf{Lower bound on $\operatorname{pred}_k$ via Lemma~\ref{lem:cauchy}.}
From \eqref{6},
\[
\operatorname{pred}_k\ge
\frac{\kappa_{\mathrm{c}}}{2}\,
\frac{\|\nabla f(x_k)\|^{2}}{L+\Delta_k^{-1}\|\nabla f(x_k)\|}.
\tag{11}
\]

\item \textbf{Derive a bound on the gradient norm.}
Combine (10) and (11).  Since $\Delta_k\ge\Delta_{\min}>0$, we have
$L+\Delta_k^{-1}\|\nabla f(x_k)\|\le L+\Delta_{\min}^{-1}\|\nabla f(x_k)\|$.
Thus,
\[
\frac{\kappa_{\mathrm{c}}}{2}\,
\frac{\|\nabla f(x_k)\|^{2}}{L+\Delta_{\min}^{-1}\|\nabla f(x_k)\|}
\;\le\;
\frac{L}{2(1-\eta_{1})}\Delta_k^{2}
\;\le\;
\frac{L}{2(1-\eta_{1})}\Delta_{\max}^{2}.
\]
Rearranging gives
\[
\|\nabla f(x_k)\|^{2}\le
\frac{L\Delta_{\max}^{2}}{\kappa_{\mathrm{c}}(1-\eta_{1})}
\Bigl(L+\Delta_{\min}^{-1}\|\nabla f(x_k)\|\Bigr).
\]
Solving the quadratic inequality in $\|\nabla f(x_k)\|$ yields
\[
\|\nabla f(x_k)\|\;\le\;
\frac{L\Delta_{\max}}{\kappa_{\mathrm{c}}(1-\eta_{1})}
+ \sqrt{\Bigl(\frac{L\Delta_{\max}}{\kappa_{\mathrm{c}}(1-\eta_{1})}\Bigr)^{2}
+\frac{L^{2}\Delta_{\max}^{2}}{\kappa_{\mathrm{c}}(1-\eta_{1})\Delta_{\min}} }.
\]
For simplicity we use the looser bound
\[
\|\nabla f(x_k)\|\;\le\;
\frac{2L\Delta_{\max}}{\kappa_{\mathrm{c}}(1-\eta_{1})}.
\tag{12}
\]

\item \textbf{Descent of the objective function.}
When a step is accepted,
\[
f(x_{k+1})=f(x_k)-\operatorname{ared}_k
\le f(x_k)-\eta_{1}\operatorname{pred}_k.
\tag{13}
\]
Using Lemma~\ref{lem:cauchy} and (12) we obtain
\[
\operatorname{pred}_k
\ge\frac{\kappa_{\mathrm{c}}}{2}
\frac{\|\nabla f(x_k)\|^{2}}{L+\Delta_{\min}^{-1}\|\nabla f(x_k)\|}
\ge
\frac{\kappa_{\mathrm{c}}}{4L}\|\nabla f(x_k)\|^{2}.
\tag{14}
\]
Substituting (14) into (13) gives
\[
f(x_{k+1})\le f(x_k)-\frac{\eta_{1}\kappa_{\mathrm{c}}}{4L}\,
\|\nabla f(x_k)\|^{2}.
\tag{15}
\]

\item \textbf{Summation over accepted iterations.}
Let $\mathcal{A}$ denote the set of indices where the step is accepted.
Summing \eqref{15} over $k\in\mathcal{A}\cap\{0,\dots,N-1\}$ yields
\[
f(x_{0})-f(x_{N})
\;\ge\;
\frac{\eta_{1}\kappa_{\mathrm{c}}}{4L}}
\sum_{k\in\mathcal{A}\cap\{0,\dots,N-1\}}
\|\nabla f(x_k)\|^{2}.
\tag{16}
\]
Since $f$ is bounded below by $f(x^{\star})$, we have
\[
\sum_{k\in\mathcal{A}\cap\{0,\dots,N-1\}}
\|\nabla f(x_k)\|^{2}
\;\le\;
\frac{4L\bigl(f(x_{0})-f(x^{\star})\bigr)}{\eta_{1}\kappa_{\mathrm{c}}}
\le\frac{4LD^{2}}{\eta_{1}\kappa_{\mathrm{c}}},
\tag{17}
\]
where the last inequality uses $f(x_{0})-f(x^{\star})\le \frac{L}{2}D^{2}$
from smoothness and bounded level set.

\item \textbf{Complexity bound.}
If $\|\nabla f(x_k)\|>\varepsilon$ for all $k<N$, then each term in the sum
\eqref{17} is at least $\varepsilon^{2}$, implying
\[
N\,\varepsilon^{2}\le
\frac{4LD^{2}}{\eta_{1}\kappa_{\mathrm{c}}}
\;\Longrightarrow\;
N\le\frac{4LD^{2}}{\eta_{1}\kappa_{\mathrm{c}}}\,\varepsilon^{-2}.
\]
Replacing $\eta_{1}$ by $1-\eta_{1}$ (the worst‑case acceptance probability) gives the bound
\eqref{5} stated in Theorem~\ref{thm:main}.  Hence the algorithm terminates
with $\|\nabla f(x_k)\|\le\varepsilon$ after at most the number of iterations
given in \eqref{5}.  Moreover, because the sum of squared gradients is finite,
$\|\nabla f(x_k)\|\to0$, establishing global convergence.
\end{enumerate}
\hfill\qed

\section*{Rate Derivation (With Constants)}
Collecting the constants that appear in \eqref{5} we have
\[
C:=\frac{2L D^{2}}{\kappa_{\mathrm{c}}(1-\eta_{1})},
\qquad
N(\varepsilon)\le C\,\varepsilon^{-2}.
\]
All quantities $L$, $D$, $\kappa_{\mathrm{c}}$, and $\eta_{1}$ are known a‑priori
or can be estimated from problem data, so the rate is explicit.

\section*{Special Cases}
\begin{itemize}
\item \textbf{Exact subproblem solution.}
If $p_k$ solves \eqref{1} exactly, then $\kappa_{\mathrm{c}}=1$.
The constant $C$ improves accordingly.
\item \textbf{Quadratic objectives.}
When $f$ is a quadratic with Hessian $Q\succeq0$, the model coincides with $f$
and $e_k(p)=0$.  Consequently $\rho_k\equiv1$, the radius grows unbounded,
and the method reduces to the classical Newton step with
global $O(\varepsilon^{-2})$ complexity.
\item \textbf{Non‑convex but $L$‑smooth.}
The proof above uses convexity only to guarantee $\operatorname{pred}_k>0$ and
the bound $f(x_{0})-f(x^{\star})\le \frac{L}{2}D^{2}$.  With a suitable
modification (e.g., requiring sufficient decrease of the model), similar
complexity results can be obtained for non‑convex smooth functions, albeit
with stationarity rather than optimality guarantees.
\end{itemize}

\end{document}